% Options for packages loaded elsewhere
\PassOptionsToPackage{unicode}{hyperref}
\PassOptionsToPackage{hyphens}{url}
%
\documentclass[
]{article}
\usepackage{lmodern}
\usepackage{amssymb,amsmath}
\usepackage{ifxetex,ifluatex}
\ifnum 0\ifxetex 1\fi\ifluatex 1\fi=0 % if pdftex
  \usepackage[T1]{fontenc}
  \usepackage[utf8]{inputenc}
  \usepackage{textcomp} % provide euro and other symbols
\else % if luatex or xetex
  \usepackage{unicode-math}
  \defaultfontfeatures{Scale=MatchLowercase}
  \defaultfontfeatures[\rmfamily]{Ligatures=TeX,Scale=1}
\fi
% Use upquote if available, for straight quotes in verbatim environments
\IfFileExists{upquote.sty}{\usepackage{upquote}}{}
\IfFileExists{microtype.sty}{% use microtype if available
  \usepackage[]{microtype}
  \UseMicrotypeSet[protrusion]{basicmath} % disable protrusion for tt fonts
}{}
\makeatletter
\@ifundefined{KOMAClassName}{% if non-KOMA class
  \IfFileExists{parskip.sty}{%
    \usepackage{parskip}
  }{% else
    \setlength{\parindent}{0pt}
    \setlength{\parskip}{6pt plus 2pt minus 1pt}}
}{% if KOMA class
  \KOMAoptions{parskip=half}}
\makeatother
\usepackage{xcolor}
\IfFileExists{xurl.sty}{\usepackage{xurl}}{} % add URL line breaks if available
\IfFileExists{bookmark.sty}{\usepackage{bookmark}}{\usepackage{hyperref}}
\hypersetup{
  hidelinks,
  pdfcreator={LaTeX via pandoc}}
\urlstyle{same} % disable monospaced font for URLs
\usepackage{longtable,booktabs}
% Correct order of tables after \paragraph or \subparagraph
\usepackage{etoolbox}
\makeatletter
\patchcmd\longtable{\par}{\if@noskipsec\mbox{}\fi\par}{}{}
\makeatother
% Allow footnotes in longtable head/foot
\IfFileExists{footnotehyper.sty}{\usepackage{footnotehyper}}{\usepackage{footnote}}
\makesavenoteenv{longtable}
\setlength{\emergencystretch}{3em} % prevent overfull lines
\providecommand{\tightlist}{%
  \setlength{\itemsep}{0pt}\setlength{\parskip}{0pt}}
\setcounter{secnumdepth}{-\maxdimen} % remove section numbering

\author{}
\date{}

\begin{document}

\hypertarget{ux91d1ux878dux6570ux636eux5206ux6790ux4e0eux667aux80fdux91cfux5316ux4ea4ux6613ux5e94-ux7528ux8bfeux7a0bux8bbeux8ba1ux8bbaux6587ux64b0ux5199ux8bf4ux660e}{%
\section{\texorpdfstring{\textbf{《金融数据分析与智能量化交易应}
\textbf{用课程设计》论文撰写说明}}{《金融数据分析与智能量化交易应 用课程设计》论文撰写说明}}\label{ux91d1ux878dux6570ux636eux5206ux6790ux4e0eux667aux80fdux91cfux5316ux4ea4ux6613ux5e94-ux7528ux8bfeux7a0bux8bbeux8ba1ux8bbaux6587ux64b0ux5199ux8bf4ux660e}}

本说明针对《金融数据分析与智能量化交易应用课程设计》论文撰写，核心要求是
\textbf{侧}

\textbf{重理论应用}
，围绕中国股市数据开展获取、分析、建模，或构建量化交易策略并分析

效果，符合学术论文基本格式规范。考虑到本科生缺乏学术论文写作训练，本说明将

对论文各组成部分进行详细拆解，明确撰写要点、内容要求和注意事项，同时结合评

分重点，帮助学生清晰掌握论文撰写逻辑和核心得分点。

论文整体要求：体裁为学术论文（Word 或 LaTex
生成均可），语言严谨、逻辑清晰、

数据真实、论证充分，所有涉及的分析、建模、策略设计需配套完整代码（建议使用

Python，可选用其他编程语言），代码需与论文内容对应，一并提交给学习委员汇总。

\hypertarget{ux4e00ux9898ux76eetitle}{%
\subsection{\texorpdfstring{\textbf{一、题目（Title）}}{一、题目（Title）}}\label{ux4e00ux9898ux76eetitle}}

\hypertarget{ux64b0ux5199ux8981ux6c42}{%
\subsubsection{\texorpdfstring{\textbf{撰写要求}}{撰写要求}}\label{ux64b0ux5199ux8981ux6c42}}

题目需 \textbf{清晰、具体、明确}
，直接反映论文的核心内容，避免模糊、宽泛或过于抽象的

表述，字数控制在 20
字以内为宜，确保读者通过题目即可知晓论文研究的核心的是

``金融股票数据分析''``量化交易策略设计''还是``股票建模应用''，且明确研究对象（如中

国 A 股某板块、具体股票）。

\hypertarget{ux5408ux683cux793aux4f8b}{%
\subsubsection{\texorpdfstring{\textbf{合格示例}}{合格示例}}\label{ux5408ux683cux793aux4f8b}}

\begin{enumerate}
\def\labelenumi{\arabic{enumi}.}
\item
  基于均线与 RSI 指标的 A 股量化交易策略设计与效果分析
\item
  中国创业板股票收益率影响因素分析及预测建模
\item
  基于 LSTM 模型的 A 股个股股价短期预测及交易策略优化
\end{enumerate}

\hypertarget{ux4e0dux5408ux683cux793aux4f8b}{%
\subsubsection{\texorpdfstring{\textbf{不合格示例}}{不合格示例}}\label{ux4e0dux5408ux683cux793aux4f8b}}

\begin{enumerate}
\def\labelenumi{\arabic{enumi}.}
\item
  金融数据分析与量化交易（过于宽泛，未明确研究内容和对象）
\item
  股票策略研究（模糊不清，未体现分析、建模或应用核心）
\end{enumerate}

\hypertarget{ux4e8cux6458ux8981abstract}{%
\subsection{\texorpdfstring{\textbf{二、摘要（Abstract）}}{二、摘要（Abstract）}}\label{ux4e8cux6458ux8981abstract}}

\hypertarget{ux64b0ux5199ux8981ux6c42-1}{%
\subsubsection{\texorpdfstring{\textbf{撰写要求}}{撰写要求}}\label{ux64b0ux5199ux8981ux6c42-1}}

摘要为论文的``浓缩版''，需 \textbf{简洁、连贯、全面}
，无需展开细节，仅需概括论文的核心内

容，让读者快速了解论文的研究目的、研究方法、核心工作和主要结论，避免冗余表

述和无关信息，不出现公式、图表和参考文献引用。

\hypertarget{ux6838ux5fc3ux5305ux542bux8981ux7d20}{%
\subsubsection{\texorpdfstring{\textbf{核心包含要素}}{核心包含要素}}\label{ux6838ux5fc3ux5305ux542bux8981ux7d20}}

\begin{enumerate}
\def\labelenumi{\arabic{enumi}.}
\tightlist
\item
  研究背景（简要 1-2
  句话，如``中国股市波动较大，量化交易凭借客观性和高效性成
\end{enumerate}

为重要交易方式''）；

\begin{enumerate}
\def\labelenumi{\arabic{enumi}.}
\setcounter{enumi}{1}
\tightlist
\item
  研究目的（明确论文要解决的问题，如``设计一种适用于 A
  股的量化交易策略，提升
\end{enumerate}

交易收益并控制风险''或``分析 A
股某板块股票数据的特征，建立股价预测模型''）；

\begin{enumerate}
\def\labelenumi{\arabic{enumi}.}
\setcounter{enumi}{2}
\tightlist
\item
  研究方法（简要说明所用的数据分析方法、建模工具或策略框架，如``基于
  Python
\end{enumerate}

获取 A
股数据，采用移动平均线指标结合机器学习模型，构建量化交易策略''）；

\begin{enumerate}
\def\labelenumi{\arabic{enumi}.}
\setcounter{enumi}{3}
\tightlist
\item
  核心结果（简要呈现实验或分析结论，如``该策略在回测中实现了
  15\%的年化收益率，
\end{enumerate}

最大回撤控制在 8\%以内，优于传统买入持有策略''）；

\begin{enumerate}
\def\labelenumi{\arabic{enumi}.}
\setcounter{enumi}{4}
\tightlist
\item
  结论（简要总结研究价值，如``该策略具有一定的实用性，可为个人投资者提供交易
\end{enumerate}

参考''）。

\hypertarget{ux6ce8ux610fux4e8bux9879}{%
\subsubsection{\texorpdfstring{\textbf{注意事项}}{注意事项}}\label{ux6ce8ux610fux4e8bux9879}}

避免过于简单（如``本文对股票数据进行了分析，设计了量化策略''），需体现具体工作

和核心结果，贴合课程``理论应用''的要求。

\hypertarget{ux4e09ux5f15ux8a00introduction}{%
\subsection{\texorpdfstring{\textbf{三、引言（Introduction）}}{三、引言（Introduction）}}\label{ux4e09ux5f15ux8a00introduction}}

引言是论文的开篇，核心是``引出研究主题''，让读者理解研究的必要性和价值，逻辑上

从``大背景''到``小问题''，再到``本文解决方案''，语言严谨且有层次感，需涵盖以下核心

内容，每个部分无需单独分标题，连贯展开即可。

\hypertarget{ux6838ux5fc3ux64b0ux5199ux5185ux5bb9}{%
\subsubsection{\texorpdfstring{\textbf{核心撰写内容}}{核心撰写内容}}\label{ux6838ux5fc3ux64b0ux5199ux5185ux5bb9}}

\begin{enumerate}
\def\labelenumi{\arabic{enumi}.}
\tightlist
\item
  研究背景：结合中国股市现状，说明研究的宏观背景和现实需求，如``随着中国资本
\end{enumerate}

市场的不断完善，A
股市场投资者结构逐渐多元化，量化交易凭借其对数据的精准分

析和纪律性交易，逐渐取代部分主观交易，成为市场关注的焦点。但当前多数量化策

略存在适应性不足、风险控制能力较弱等问题，尤其是针对 A
股中小盘股票的策略较

为匮乏，因此，开展股票数据的分析、建模与量化策略设计具有重要的现实意义''；

\begin{enumerate}
\def\labelenumi{\arabic{enumi}.}
\setcounter{enumi}{1}
\tightlist
\item
  研究意义：分为理论意义和实践意义，贴合本科生水平，无需过于深入，重点突出
\end{enumerate}

``应用价值''（符合课程侧重理论应用的要求）：

（1）理论意义：简要说明研究对现有金融数据分析、量化交易理论的补充，如``丰富

了 A
股市场量化交易策略的应用场景，验证了某类机器学习模型在股票数据预测中的

有效性，为后续相关研究提供了基础参考''；

（2）实践意义：说明研究的实际应用价值，如``为个人投资者提供一种可行的量化交

易思路，帮助投资者理性分析股票数据、控制交易风险，提升投资收益''；

\begin{enumerate}
\def\labelenumi{\arabic{enumi}.}
\setcounter{enumi}{2}
\tightlist
\item
  前人工作简要说明：无需展开，简要概括现有相关研究的核心方向，如``目前国内外
\end{enumerate}

学者围绕量化交易策略开展了大量研究，主要集中在均线策略、动量策略、反转策略

以及结合机器学习的复合策略，但现有策略在 A
股市场的适应性仍有提升空间，部分

策略未充分考虑中国股市的政策影响和波动特征''；

\begin{enumerate}
\def\labelenumi{\arabic{enumi}.}
\setcounter{enumi}{3}
\tightlist
\item
  本论文的理论内容：明确本文的研究核心、所用的理论和方法，如``本文以中国
  A 股
\end{enumerate}

沪深 300 成分股为研究对象，基于金融数据分析理论和量化交易原理，采用
Python 获

取股票历史交易数据，通过数据清洗、特征分析，构建结合均线指标和 LSTM
模型的

量化交易策略，对策略的回测效果进行分析''；

\begin{enumerate}
\def\labelenumi{\arabic{enumi}.}
\setcounter{enumi}{4}
\tightlist
\item
  本文创新点：结合本科生水平，创新点无需过于前沿，重点突出``差异化''和``实用性''，
\end{enumerate}

避免照搬现有研究，如``（1）优化了传统均线策略的参数设置，结合机器学习模型提

升了策略的预测准确性；（2）针对 A
股市场波动特征，加入风险控制模块，降低策略

的最大回撤；（3）选用沪深 300
成分股为研究对象，确保策略的实用性和可复制性''；

\begin{enumerate}
\def\labelenumi{\arabic{enumi}.}
\setcounter{enumi}{5}
\tightlist
\item
  论文结构安排（可选）：简要说明论文各章节的核心内容，帮助读者理清论文逻辑，
\end{enumerate}

如``本文后续章节安排如下：第二章总结相关工作；第三章介绍数据获取与清洗过程；

第四章详细阐述本文提出的量化交易策略和建模方法；第五章展示实验结果与对比分

析；第六章总结全文并展望未来研究方向''。

\hypertarget{ux56dbux76f8ux5173ux5de5ux4f5crelated-work}{%
\subsection{\texorpdfstring{\textbf{四、相关工作（Related
Work）}}{四、相关工作（Related Work）}}\label{ux56dbux76f8ux5173ux5de5ux4f5crelated-work}}

相关工作是对``前人研究''的详细总结和分析，核心是``梳理现有研究、对比优劣、引出

本文研究的必要性''，体现学生的文献阅读和总结能力，需结合课程主题（金融股票数

据分析、量化交易策略），避免偏离主题。

\hypertarget{ux64b0ux5199ux8981ux70b9}{%
\subsubsection{\texorpdfstring{\textbf{撰写要点}}{撰写要点}}\label{ux64b0ux5199ux8981ux70b9}}

\begin{enumerate}
\def\labelenumi{\arabic{enumi}.}
\tightlist
\item
  文献来源：建议阅读中外文学术论文（可通过知网、Google Scholar、IEEE
  Xplore
\end{enumerate}

等平台获取），优先选择近 5 年的相关研究（保证研究的时效性），至少参考 5
篇核心

文献，无需过多（本科生无需追求文献数量，重点在于总结质量）；

\begin{enumerate}
\def\labelenumi{\arabic{enumi}.}
\setcounter{enumi}{1}
\tightlist
\item
  总结内容：按``研究方向''分类总结前人工作，避免零散罗列，如可分为``传统量化交
\end{enumerate}

易策略研究''``机器学习在股票数据分析中的应用研究''``A
股市场量化策略研究''三类，每

类总结 2-3 篇文献的核心观点、所用方法和研究结论；

\begin{enumerate}
\def\labelenumi{\arabic{enumi}.}
\setcounter{enumi}{2}
\tightlist
\item
  优缺点分析（核心）：针对每类前人工作，客观分析其优点和不足，重点突出``不足''，
\end{enumerate}

为本文的创新点和研究方向提供支撑，贴合课程``理论应用''的要求，不足部分需与本文

研究内容相关，如：

（1）传统均线策略：优点是简单易懂、可操作性强，适用于短期交易；不足是参数设

置固定，适应性差，难以应对 A 股市场的突发波动，收益稳定性不足；

（2）基于机器学习的量化策略：优点是能充分挖掘股票数据的非线性特征，预测准确

性较高；不足是模型复杂，需要大量数据支撑，部分模型过度拟合，实际应用效果不

佳，且未充分结合中国股市的政策因素；

\begin{enumerate}
\def\labelenumi{\arabic{enumi}.}
\setcounter{enumi}{3}
\tightlist
\item
  本文与前人工作的关联：简要说明本文如何弥补前人工作的不足，如``本文针对前人
\end{enumerate}

策略适应性差、风险控制不足的问题，优化策略参数并加入风险控制模块，结合传统

技术指标和机器学习模型，提升策略在 A 股市场的实用性和收益稳定性''。

\hypertarget{ux6ce8ux610fux4e8bux9879-1}{%
\subsubsection{\texorpdfstring{\textbf{注意事项}}{注意事项}}\label{ux6ce8ux610fux4e8bux9879-1}}

避免单纯罗列文献（如``文献{[}1{]}研究了均线策略，文献{[}2{]}研究了 LSTM
模型的应用''），

需对文献内容进行提炼和分析，优缺点分析需客观，不夸大不足，也不忽视优点，贴

合本科生的文献总结能力。

\hypertarget{ux4e94ux6570ux636eux83b7ux53d6ux4e0eux6e05ux6d17data-acquisition-and-cleaning}{%
\subsection{\texorpdfstring{\textbf{五、数据获取与清洗（Data Acquisition
and}
\textbf{Cleaning）}}{五、数据获取与清洗（Data Acquisition and Cleaning）}}\label{ux4e94ux6570ux636eux83b7ux53d6ux4e0eux6e05ux6d17data-acquisition-and-cleaning}}

本部分是论文的``基础''，核心是说明研究所用数据的来源、类型、处理过程，体现学生

的数据获取和预处理能力，贴合课程``金融股票数据获取、分析''的核心要求，需详细、

具体，可结合图表辅助说明（如数据描述性统计表格）。

\hypertarget{ux6838ux5fc3ux64b0ux5199ux5185ux5bb9-1}{%
\subsubsection{\texorpdfstring{\textbf{核心撰写内容}}{核心撰写内容}}\label{ux6838ux5fc3ux64b0ux5199ux5185ux5bb9-1}}

\begin{enumerate}
\def\labelenumi{\arabic{enumi}.}
\tightlist
\item
  数据获取说明：
\end{enumerate}

（1）数据来源：明确数据的获取渠道（需可实现，贴合本科生能力），如``本文数据来

源于 Tushare 金融数据平台、聚宽量化平台或 Wind
金融终端（选用免费可获取的渠

道，方便其他读者复现），获取的是中国 A 股沪深 300
成分股的历史交易数据''；

（2）数据类型：明确数据的具体内容，如``数据类型包括日度交易数据，具体字段有

日期、开盘价、收盘价、最高价、最低价、成交量、成交额、换手率等，同时可补充

股票的基本信息（如行业分类）''；

（3）数据时间范围：明确数据的起止时间，如``数据时间范围为 2020 年 1 月 1
日至

2024 年 12 月 31 日，共 1200
个交易日左右，确保数据量充足，满足建模和回测需求''；

（4）获取方法：简要说明获取数据的技术手段，结合课程要求和评分重点（代码提

交），如``采用 Python 编程语言，通过 Tushare API
调用接口获取数据，编写代码读取

数据并保存为 CSV 格式，方便后续处理，具体代码见附录''；

\begin{enumerate}
\def\labelenumi{\arabic{enumi}.}
\setcounter{enumi}{1}
\tightlist
\item
  数据初步分析：对获取的原始数据进行简要描述性统计，说明数据的基本特征，如
\end{enumerate}

``本次共获取沪深 300 成分股中 50 只股票的日度数据，总数据量为 60000
条，通过描

述性统计可知，样本股票的平均日收益率为 0.08\%，成交量均值为 5000
万元，换手

率均值为
2.3\%，部分股票存在数据缺失和异常值（如涨跌停导致的价格异常）''；

\begin{enumerate}
\def\labelenumi{\arabic{enumi}.}
\setcounter{enumi}{2}
\tightlist
\item
  数据清洗过程（核心）：详细说明数据预处理的步骤、方法和目的，清洗后的效果，
\end{enumerate}

这是后续建模和策略设计的基础，需具体可操作，如：

（1）缺失值处理：说明缺失值的存在情况和处理方法，如``部分股票因停牌导致数据

缺失，共缺失 320 条数据，采用`前向填充法'（forward
fill）填充缺失值，确保数据的

连续性，避免因缺失值影响建模效果''；

（2）异常值处理：识别并处理异常数据（如涨跌停异常、成交量异常），如``采用`3σ

原则'识别异常值，对超出均值±3
倍标准差的成交量、换手率数据，采用中位数替换法

处理，共处理异常值 180 条，确保数据的合理性''；

（3）数据标准化/归一化（可选）：若后续建模需要，说明数据标准化的方法和目的，

如``为消除数据量纲影响（如收盘价和成交量的量纲差异），采用 Min-Max
归一化方法，

将数据映射到{[}0,1{]}区间，提升模型的收敛速度和预测准确性''；

（4）数据筛选与重构（可选）：根据研究目的，筛选有用的数据字段或重构特征，如

``筛选出收盘价、成交量、换手率 3 个核心字段，构建`日收益率'\,`5
日均线'等衍生特征，

用于后续策略设计和建模''；

\begin{enumerate}
\def\labelenumi{\arabic{enumi}.}
\setcounter{enumi}{3}
\tightlist
\item
  数据清洗结果：说明清洗后的数据情况，如``数据清洗后，总数据量为 59800
  条，无
\end{enumerate}

缺失值、异常值，数据连续性和合理性得到保证，清洗后的数据用于后续的特征分析、

建模和策略回测，清洗过程的相关代码见附录''。

\hypertarget{ux6ce8ux610fux4e8bux9879-2}{%
\subsubsection{\texorpdfstring{\textbf{注意事项}}{注意事项}}\label{ux6ce8ux610fux4e8bux9879-2}}

需结合代码，说明每一步清洗操作的目的和方法，避免只说``进行了数据清洗''而不说明

具体过程；数据来源需真实可获取，避免虚构数据（评分重点之一）；可结合表格展示

描述性统计结果，提升可读性。

\hypertarget{ux516dux65b9ux6cd5method}{%
\subsection{\texorpdfstring{\textbf{六、方法（Method）}}{六、方法（Method）}}\label{ux516dux65b9ux6cd5method}}

本部分是论文的``核心''，体现学生的理论应用和建模能力，需详细、清晰地介绍本文提

出的数据分析方法、建模过程或量化交易策略，包括核心思路、数学公式、模型图、

创新点等，贴合课程``建模与应用''的要求，字数控制在 1000-1500
字，确保读者能够

根据本部分内容复现研究过程。

\hypertarget{ux64b0ux5199ux8981ux70b9-1}{%
\subsubsection{\texorpdfstring{\textbf{撰写要点}}{撰写要点}}\label{ux64b0ux5199ux8981ux70b9-1}}

\begin{enumerate}
\def\labelenumi{\arabic{enumi}.}
\tightlist
\item
  方法概述：首先简要说明本文所用方法的核心思路和目的，如``本文提出一种结合均
\end{enumerate}

线指标和 LSTM
模型的量化交易策略，核心思路是：通过均线指标识别股票的趋势方

向，利用 LSTM
模型预测股票的短期收盘价，结合两者生成交易信号（买入/卖出），

同时加入风险控制模块，控制交易风险，提升策略收益''；

\begin{enumerate}
\def\labelenumi{\arabic{enumi}.}
\setcounter{enumi}{1}
\tightlist
\item
  相关理论基础（简要）：简要介绍方法所用的理论、指标或模型的基础，无需过于深
\end{enumerate}

入，贴合本科生水平，如``（1）移动平均线（MA）：用于识别股票价格趋势，本文选

用 5 日均线（MA5）和 20 日均线（MA20），MA5 反映短期趋势，MA20
反映中期趋

势；（2）LSTM
模型：长短期记忆网络，适用于时间序列数据预测，能够捕捉股票价

格的长期依赖关系，避免梯度消失问题''；

\begin{enumerate}
\def\labelenumi{\arabic{enumi}.}
\setcounter{enumi}{2}
\tightlist
\item
  详细方法描述（核心）：
\end{enumerate}

（1）量化交易策略设计（若研究策略）：分步骤介绍策略的框架，包括信号生成、买

入/卖出规则、仓位管理、风险控制等，结合数学公式和流程图说明，如：

① 趋势识别（均线指标）：定义 5 日均线和 20 日均线的计算公式：

{[}1{]} 4 𝐶𝑙𝑜𝑠𝑒(𝑡−𝑖)

5 {[}∑{]} 𝑖=0

𝑀𝐴5(𝑡) = {[}1{]}

𝑖=0

{[}1{]} 19 𝐶𝑙𝑜𝑠𝑒(𝑡−𝑖)

20 {[}∑{]} 𝑖=0

𝑀𝐴20(𝑡) = {[}1{]}

𝑖=0

其中， 𝑀𝐴5(𝑡) 为 t 时刻的 5 日均线值， 𝑀𝐴20(𝑡) 为 t 时刻的 20 日均线值，
𝐶𝑙𝑜𝑠𝑒(𝑡−

𝑖) 为 t-i 时刻的股票收盘价。

趋势判断规则：当 𝑀𝐴5(𝑡) \textgreater{} 𝑀𝐴20(𝑡) 时，判断为上升趋势；当
𝑀𝐴5(𝑡) \textless{} 𝑀𝐴20(𝑡) 时，

判断为下降趋势。

② 价格预测（LSTM
模型）：简要介绍模型的结构（输入层、隐藏层、输出层），输入

特征（如前 5
日的收盘价、成交量），输出（下一日的收盘价预测值），结合模型流程

图说明（可手绘或用 Visio、Mermaid 绘制，插入论文中），如``LSTM
模型输入层为 5

个特征（前 5 日收盘价、前 5 日成交量），隐藏层包含 128
个神经元，输出层为 1 个预

测值（下一日收盘价），模型采用 Adam
优化器，损失函数为均方误差（MSE），计算

公式为：

𝑛 𝑛 {[}1{]} {[}∑{]} 𝑖=1(𝑦𝑖 −𝑦̂𝑖) {[}2{]}

𝑀𝑆𝐸= {[}1{]}

𝑖=1

其中， 𝑦𝑖 为实际收盘价， 𝑦̂𝑖 为预测收盘价，n 为样本数量。''

③
交易信号生成：结合趋势识别和价格预测结果，制定买入/卖出规则，如``买入信号：

MA5(t) \textgreater{} MA20(t)（上升趋势）且预测收盘价 \textgreater{}
当日收盘价（价格上涨预期）；卖出信

号：MA5(t) \textless{} MA20(t)（下降趋势）或预测收盘价 \textless{}
当日收盘价（价格下跌预期），同

时设置止损线（亏损达到 5\%时强制卖出）和止盈线（盈利达到
10\%时强制卖出）。''

④ 仓位管理：明确仓位设置，如``采用固定仓位策略，每次交易投入总资金的
20\%，

避免单一股票仓位过高导致的风险。''

（2）数据分析与建模（若不做策略，仅做分析和预测）：详细介绍数据特征分析方法、

建模过程，如``采用相关性分析（计算各特征与收益率的相关系数）、时序分析（分析

股票价格的趋势和周期性），构建多元线性回归模型预测股票收益率，模型公式为：

𝑦= 𝛽0 + 𝛽1𝑥1 + 𝛽2𝑥2+. . . +𝛽𝑘𝑥𝑘 + 𝜀

其中，y 为股票收益率， 𝑥1, 𝑥2, . . ., 𝑥𝑘
为特征变量（如成交量、换手率）， 𝛽0 为截距项，

𝛽1, . . ., 𝛽𝑘 为回归系数， 𝜀
为随机误差项，详细说明模型的训练过程、参数估计方法。''

\begin{enumerate}
\def\labelenumi{\arabic{enumi}.}
\setcounter{enumi}{3}
\tightlist
\item
  创新点详细说明（核心）：结合本部分的方法描述，详细阐述本文的创新点，比引言
\end{enumerate}

中的创新点更具体，如``本文的创新点主要体现在两个方面：（1）策略创新：将传统均

线指标（趋势识别）与 LSTM
模型（价格预测）相结合，弥补了传统均线策略无法预

测价格短期波动的不足，同时避免了单一机器学习模型过度拟合的问题；（2）风险控

制创新：加入止损止盈线和固定仓位管理，针对 A
股市场波动较大的特点，有效降低

了策略的最大回撤，提升了收益稳定性，这是前人部分策略所缺失的。''

\hypertarget{ux6ce8ux610fux4e8bux9879-3}{%
\subsubsection{\texorpdfstring{\textbf{注意事项}}{注意事项}}\label{ux6ce8ux610fux4e8bux9879-3}}

\begin{enumerate}
\def\labelenumi{\arabic{enumi}.}
\item
  数学公式需规范（例如用 LaTex
  编写），明确每个符号的含义，避免公式错误；
\item
  模型图、策略流程图需清晰，标注明确，辅助说明方法原理；
\item
  方法描述需详细可复现，结合代码思路，为后续实验结果和代码提交做铺垫（评分
\end{enumerate}

重点之一）；

\begin{enumerate}
\def\labelenumi{\arabic{enumi}.}
\setcounter{enumi}{3}
\tightlist
\item
  避免方法描述过于笼统，如``采用 LSTM
  模型进行预测''，需说明模型结构、输入输
\end{enumerate}

出、参数设置等细节。

\hypertarget{ux4e03ux5b9eux9a8cux7ed3ux679cexperimental-results}{%
\subsection{\texorpdfstring{\textbf{七、实验结果（Experimental
Results）}}{七、实验结果（Experimental Results）}}\label{ux4e03ux5b9eux9a8cux7ed3ux679cexperimental-results}}

本部分是论文的``关键''，核心是展示研究的实际效果，体现学生的实验设计和结果分析

能力，贴合课程``理论应用''和评分重点（策略效果越好，评分越高），需详细、客观地

描述实验过程、结果，进行对比分析，可结合表格、图表（如收益率曲线、回测结果

对比表）辅助说明，增强说服力。

\hypertarget{ux64b0ux5199ux8981ux70b9-2}{%
\subsubsection{\texorpdfstring{\textbf{撰写要点}}{撰写要点}}\label{ux64b0ux5199ux8981ux70b9-2}}

\begin{enumerate}
\def\labelenumi{\arabic{enumi}.}
\tightlist
\item
  实验环境说明：明确实验所用的硬件、软件、编程语言、相关库，方便复现，如``实
\end{enumerate}

验硬件：CPU 为 Intel Core i7，内存 16GB；实验软件：Python
3.9，相关库包括

Pandas （数据处理）、 Matplotlib （可视化）、 Scikit-learn （建模）、
TensorFlow

（LSTM 模型）；实验数据：采用数据清洗后的沪深 300
成分股日度数据（2020-2024

年），其中 80\%用于模型训练，20\%用于测试和策略回测''；

\begin{enumerate}
\def\labelenumi{\arabic{enumi}.}
\setcounter{enumi}{1}
\tightlist
\item
  实验设计：明确实验的目的、指标、对比方法，如``实验目的：验证本文提出的量化
\end{enumerate}

交易策略的有效性和优越性；实验指标：选用年化收益率、最大回撤、夏普比率、胜

率 4 个核心指标（贴合量化交易策略效果评价标准），各指标含义如下：

（ 1 ）年化收益率：衡量策略的收益水平，计算公式为：年化收益率 = (1 +

250

总收益率 ) 交易天数 −1 （250 为 A 股年平均交易日）；

（2）最大回撤：衡量策略的风险水平，指策略在回测期间最大的亏损幅度，数值越小，

风险越低；

（ 3 ）夏普比率：衡量策略的风险调整后收益，计算公式为：夏普比率 =

策略年化收益率 - 无风险利率

（无风险利率选用 3 年期国债收益率，取 2.75\%），数值越大，
策略年化波动率

策略性价比越高；

（4）胜率：衡量策略的交易准确性，计算公式为：胜率 =

盈利交易次数

总交易次数 {[}× 100\%{]} {[}； {]}

对比方法：选用 2
种常见的量化策略作为对比，分别是传统均线策略（MA5+MA20）、

买入持有策略（买入后长期持有，不进行交易），确保对比的合理性和针对性''；

\begin{enumerate}
\def\labelenumi{\arabic{enumi}.}
\setcounter{enumi}{2}
\tightlist
\item
  实验过程简要描述：简要说明实验的步骤，如``（1）将清洗后的数据分为训练集和
\end{enumerate}

测试集，训练集用于 LSTM
模型训练和策略参数优化，测试集用于策略回测；（2）优

化模型和策略参数（如 LSTM
模型的神经元数量、均线指标的参数、止损止盈比例），

确保策略效果最优；（3）分别运行本文策略、传统均线策略、买入持有策略，记录各

策略的回测指标；（4）对实验结果进行统计、可视化，对比分析各策略的优劣''；

\begin{enumerate}
\def\labelenumi{\arabic{enumi}.}
\setcounter{enumi}{3}
\tightlist
\item
  实验结果详细展示：
\end{enumerate}

（1）数值结果：用表格展示各策略的核心指标，如：

\begin{longtable}[]{@{}lllll@{}}
\toprule
策略类型 & 年化收益率（\%） & 最大回撤（\%） & 夏普比率 &
胜率（\%）\tabularnewline
\midrule
\endhead
本文提出的策略 & 15.2 & 7.8 & 1.86 & 58.3\tabularnewline
传统均线策略 & 10.5 & 12.3 & 1.23 & 52.1\tabularnewline
买入持有策略 & 8.7 & 15.6 & 1.02 & 50.5\tabularnewline
\bottomrule
\end{longtable}

（2）可视化结果：说明所用的可视化图表及核心信息，如``绘制各策略的累计收益率

曲线（如图 1
所示），可以看出，本文提出的策略累计收益率始终高于传统均线策略和

买入持有策略，在市场下跌期间，本文策略的回撤幅度更小，体现了更好的风险控制

能力；绘制 LSTM 模型的预测值与实际值对比图（如图 2
所示），可以看出，模型预测

值与实际值拟合度较高，平均预测误差为 2.3\%，验证了模型的有效性''；

\begin{enumerate}
\def\labelenumi{\arabic{enumi}.}
\setcounter{enumi}{4}
\tightlist
\item
  结果分析与对比（核心）：
\end{enumerate}

（1）自身结果分析：分析本文策略或模型的效果，如``从实验结果可以看出，本文提

出的量化交易策略年化收益率达到 15.2\%，最大回撤控制在 7.8\%，夏普比率为
1.86，

胜率为
58.3\%，说明该策略能够在保证收益的同时，有效控制风险，具有较好的实用

性；LSTM 模型的平均预测误差为
2.3\%，能够较为准确地预测股票短期收盘价，为交

易信号生成提供了可靠支撑''；

（2）与对比方法的对比分析：突出本文方法的优越性，结合相关工作中前人工作的不

足，说明本文方法如何弥补这些不足，如``与传统均线策略相比，本文策略的年化收益

率提升了 4.7 个百分点，最大回撤降低了 4.5 个百分点，原因是本文策略结合了
LSTM

模型的价格预测能力，优化了交易信号生成时机，同时加入了止损止盈模块，降低了

风险；与买入持有策略相比，本文策略的年化收益率提升了 6.5
个百分点，最大回撤

降低了 7.8
个百分点，体现了量化交易策略的纪律性和优势，避免了主观交易的盲目

性''；

（3）异常结果分析（可选）：若实验中出现异常结果（如某段时间策略收益大幅下跌），

简要分析原因，如``在 2022 年 4 月-5
月期间，本文策略出现小幅亏损，主要原因是当

时 A
股市场受疫情影响出现短期大幅波动，超出了模型的预测范围，后续可通过加入

政策因素特征，优化模型的适应性''。

\hypertarget{ux6ce8ux610fux4e8bux9879-4}{%
\subsubsection{\texorpdfstring{\textbf{注意事项}}{注意事项}}\label{ux6ce8ux610fux4e8bux9879-4}}

\begin{enumerate}
\def\labelenumi{\arabic{enumi}.}
\tightlist
\item
  实验结果需真实、客观，避免虚构数据（评分重点之一），所有结果需能通过提交的
\end{enumerate}

代码复现；

\begin{enumerate}
\def\labelenumi{\arabic{enumi}.}
\setcounter{enumi}{1}
\item
  对比方法需合理，贴合研究主题，避免选用无关的方法进行对比；
\item
  指标解释需清晰，结果分析需深入，避免只展示结果而不分析原因；
\item
  图表需清晰、规范，标注图表标题、坐标轴含义，插入论文对应位置，辅助说明结
\end{enumerate}

果。

\hypertarget{ux516bux603bux7ed3conclusion}{%
\subsection{\texorpdfstring{\textbf{八、总结（Conclusion）}}{八、总结（Conclusion）}}\label{ux516bux603bux7ed3conclusion}}

本部分是论文的``收尾''，核心是对全文内容进行总结，客观分析研究的优点和不足，对

未来研究进行展望，逻辑连贯，避免重复前文内容，体现学生的总结和反思能力。

\hypertarget{ux6838ux5fc3ux64b0ux5199ux5185ux5bb9-2}{%
\subsubsection{\texorpdfstring{\textbf{核心撰写内容}}{核心撰写内容}}\label{ux6838ux5fc3ux64b0ux5199ux5185ux5bb9-2}}

\begin{enumerate}
\def\labelenumi{\arabic{enumi}.}
\tightlist
\item
  全文总结：简要回顾论文的核心工作、研究方法和主要结论，与引言、方法、实验
\end{enumerate}

结果相呼应，如``本文围绕中国 A
股数据的分析、建模与量化交易策略设计展开研究，

以沪深 300 成分股为研究对象，通过 Python
获取股票历史交易数据，完成数据清洗和

初步分析，构建了结合均线指标和 LSTM
模型的量化交易策略，通过实验验证了策略

的有效性。实验结果表明，该策略的年化收益率达到 15.2\%，最大回撤控制在
7.8\%，

优于传统均线策略和买入持有策略，能够为个人投资者提供可行的交易参考''；

\begin{enumerate}
\def\labelenumi{\arabic{enumi}.}
\setcounter{enumi}{1}
\tightlist
\item
  研究优点（客观总结）：结合本文的创新点和实验结果，总结研究的优势，如``本文
\end{enumerate}

的优点主要体现在：（1）研究贴合中国 A 股市场实际，选用沪深 300
成分股为研究对

象，策略实用性强、可复制性高；（2）方法上结合传统技术指标和机器学习模型，弥

补了单一方法的不足，同时加入风险控制模块，提升了策略的稳定性；（3）实验设计

合理，通过与常见策略对比，充分验证了本文策略的优越性，数据真实、过程可复现''；

\begin{enumerate}
\def\labelenumi{\arabic{enumi}.}
\setcounter{enumi}{2}
\tightlist
\item
  研究不足（客观反思，重点）：结合本科生的研究水平，客观分析研究的局限性，避
\end{enumerate}

免回避不足，不足需具体、合理，如``本文的研究仍存在一些不足：（1）数据范围较窄，

仅选用沪深 300
成分股，未考虑中小盘股票，策略的适用性有待进一步验证；（2）模

型和策略参数优化不够全面，仅优化了部分核心参数，可能存在更优的参数组合；（3）

未考虑政策、宏观经济等外部因素对股票价格的影响，模型的抗干扰能力有待提升；

（4）策略仅适用于短期交易，长期交易效果未进行验证''；

\begin{enumerate}
\def\labelenumi{\arabic{enumi}.}
\setcounter{enumi}{3}
\tightlist
\item
  未来展望（结合不足，提出改进方向）：针对研究不足，提出具体、可行的未来改进
\end{enumerate}

思路，贴合课程主题，如``针对本文的不足，未来可从以下方面开展进一步研究：（1）

扩大数据范围，加入中小盘股票数据，验证策略在不同类型股票中的适用性；（2）优

化模型和策略参数，采用网格搜索、遗传算法等方法进行全局参数优化，提升策略效

果；（3）引入政策、宏观经济等外部特征，完善模型输入，提升模型的抗干扰能力；

（4）拓展策略的应用场景，研究策略在长期交易、多品种交易中的效果，进一步完善

策略框架''。

\hypertarget{ux4e5dux8bc4ux5206ux91cdux70b9ux8bf4ux660e}{%
\subsection{\texorpdfstring{\textbf{九、评分重点说明}}{九、评分重点说明}}\label{ux4e5dux8bc4ux5206ux91cdux70b9ux8bf4ux660e}}

本课程论文评分核心侧重``理论应用''和``实践效果''，结合本科生水平，重点关注以下
7

点，直接影响论文成绩：

\begin{enumerate}
\def\labelenumi{\arabic{enumi}.}
\tightlist
\item
  主题相关性：论文需围绕``金融股票数据分析、建模与应用''展开（如量化交易策略
\end{enumerate}

设计、股票价格预测、收益率影响因素分析等），优先选用中国股市数据，偏离主题将

大幅扣分；

\begin{enumerate}
\def\labelenumi{\arabic{enumi}.}
\setcounter{enumi}{1}
\tightlist
\item
  格式规范性：严格按照本文说明的论文结构（题目、摘要、引言等 8
  个部分）撰写，
\end{enumerate}

符合学术论文格式要求，Word 或 LaTex 生成均可，格式混乱、结构缺失将扣分；

\begin{enumerate}
\def\labelenumi{\arabic{enumi}.}
\setcounter{enumi}{2}
\tightlist
\item
  代码提交：必须提交与论文内容对应的完整代码（建议
  Python），代码需可运行、
\end{enumerate}

可复现实验结果；

\begin{enumerate}
\def\labelenumi{\arabic{enumi}.}
\setcounter{enumi}{3}
\tightlist
\item
  策略/模型效果：量化交易策略的回测效果（年化收益率、最大回撤等指标）是重要
\end{enumerate}

评分标准，效果越好，得分越高；若仅做数据分析与建模，需保证模型预测准确性、

分析合理性；

\begin{enumerate}
\def\labelenumi{\arabic{enumi}.}
\setcounter{enumi}{4}
\tightlist
\item
  数据真实性：数据来源需真实可获取，数据清洗过程详细、合理，禁止虚构数据，
\end{enumerate}

数据造假将直接判为不合格；

\begin{enumerate}
\def\labelenumi{\arabic{enumi}.}
\setcounter{enumi}{5}
\tightlist
\item
  创新性：无需过于前沿，重点体现差异化和实用性，避免照搬现有研究，创新点明
\end{enumerate}

确、合理可得分；

\begin{enumerate}
\def\labelenumi{\arabic{enumi}.}
\setcounter{enumi}{6}
\tightlist
\item
  逻辑与表达：论文逻辑清晰、语言严谨，无语法错误，各部分内容连贯，能够清晰
\end{enumerate}

阐述研究工作，避免冗余、模糊的表述。

\hypertarget{ux5341ux5176ux4ed6ux6ce8ux610fux4e8bux9879}{%
\subsection{\texorpdfstring{\textbf{十、其他注意事项}}{十、其他注意事项}}\label{ux5341ux5176ux4ed6ux6ce8ux610fux4e8bux9879}}

\begin{enumerate}
\def\labelenumi{\arabic{enumi}.}
\tightlist
\item
  参考文献：至少参考 3-5 篇近 5
  年的相关学术论文，在引言、相关工作部分合理引
\end{enumerate}

用，引用格式规范（可参考知网文献引用格式）；

\begin{enumerate}
\def\labelenumi{\arabic{enumi}.}
\setcounter{enumi}{1}
\tightlist
\item
  代码提交：代码可单独提交（命名格式为``姓名+学号+论文题目.py''），或作为附录
\end{enumerate}

插入论文末尾；

\begin{enumerate}
\def\labelenumi{\arabic{enumi}.}
\setcounter{enumi}{2}
\tightlist
\item
  提交方式：论文（Word/LaTex 生成的 PDF
  文件）和代码一并提交给学习委员，由
\end{enumerate}

学习委员汇总后提交；

\begin{enumerate}
\def\labelenumi{\arabic{enumi}.}
\setcounter{enumi}{3}
\tightlist
\item
  原创性要求：禁止抄袭、剽窃他人研究成果，抄袭将直接判为不合格，需独立完成
\end{enumerate}

论文撰写和代码编写，可参考现有文献，但需注明引用。

\end{document}
